
\documentclass[a4paper,10pt]{report}

%\usepackage[landscape]{geometry}
\usepackage[cm-default]{fontspec}
\setromanfont{FreeSerif}
\setsansfont{FreeSans}
\setmonofont{FreeMono}
\usepackage[utf8x]{inputenc}
\usepackage{fontspec}
\usepackage{xunicode}
\usepackage{xltxtra}
%\usepackage{xgreek}
\usepackage{amsmath}
\usepackage{unicode-math}
\usepackage{ulem}
\usepackage{color}
\usepackage{verbatim}
\usepackage{nopageno}
\usepackage{graphicx}
\usepackage{textpos}
\setlength{\TPHorizModule}{1cm}
\setlength{\TPVertModule}{1cm}
%\usepackage[colorgrid,texcoord]{eso-pic}
\usepackage[outline]{contour}
\usepackage{wrapfig}
\usepackage{url}
\usepackage{color}
\usepackage{tikz}
\usepackage{fancybox,fancyhdr}
\usepackage{subfigure}
\usepackage{pstricks}
\usepackage{epsfig}
\usepackage{multicol}
\usepackage{listings}
\usepackage{enumerate}
\usepackage{hyperref}
\hypersetup{
  bookmarks=true,
  bookmarksopen=true,
  pdfborder=false,
  pdfpagemode=UseNone,
  raiselinks=true,
  pdfhighlight={/P},
  colorlinks,
  citecolor=black,
  filecolor=black,
  linkcolor=black,
  urlcolor=black
}


\usepackage{fontspec}
\usepackage{xunicode}
\usepackage{xltxtra}


% Margins
%\setlength{\textwidth}{24cm}

%\setlength{\voffset}{0in}
%\setlength{\textheight}{6.8in}
% Colors
\definecolor{cbrown}{rgb}{0.49,0.24,0.07}
\definecolor{cmpez}{rgb}{0.92,0.88,0.79}
\definecolor{cmpez2}{rgb}{0.62,0.33,0.34}
\definecolor{cwhite}{rgb}{1,1,1}
\definecolor{cblack}{rgb}{0,0,0}
\definecolor{cred}{rgb}{0.9,0.15,0.15}
\definecolor{clblue}{rgb}{0.57,0.85,0.97}
\definecolor{corange}{rgb}{0.86,0.49,0.18}
\definecolor{clorange}{rgb}{0.95,0.84,0.61}
\definecolor{cyellow}{rgb}{1,0.95,0.16}
\definecolor{cgreen}{rgb}{0.69,0.8,0.31}
\definecolor{clbrown}{rgb}{0.78,0.67,0.43}
\definecolor{cmagenta}{rgb}{0.79,0.05,0.54}
\definecolor{cgray}{rgb}{0.76,0.73,0.63}
\definecolor{bluesite}{rgb}{0.05,0.43,0.67}
\definecolor{redsite}{rgb}{0.90,0.15,0.15}
\definecolor{backsite}{rgb}{0.07,0.07,0.07}
% ----------- from Costas -----------------------------

%---------------   Main Fonts    -----------------%   

% \setmainfont[Mapping=tex-text]{Calibri}
% \setmainfont[Mapping=tex-text]{Times New Roman}
%\setmainfont[Mapping=tex-text]{Droid Serif}
%\setmainfont[Mapping=tex-text]{Cambria}
%\setmainfont[Mapping=tex-text]{cm-unicode}
% \setmainfont[Mapping=tex-text]{Gentium}
% \setmainfont[Mapping=tex-text]{GFS Didot}
% \setmainfont[Mapping=tex-text]{Comic Sans MS}
 \setmainfont[Mapping=tex-text]{Ubuntu}
% \setmainfont[Mapping=tex-text]{Myriad Pro}
% \setmainfont[Mapping=tex-text]{CMU Concrete}
%\setmainfont[Mapping=tex-text]{DejaVu Sans}
%\setmainfont[Mapping=tex-text]{KerkisSans}
%\setmainfont[Mapping=tex-text]{KerkisCaligraphic}
% \setmainfont[Mapping=tex-text]{Segoe Print}
% \setmainfont[Mapping=tex-text]{Gabriola}

%------------------------------------------------%

% ------------ Mathematics Fonts ----------------%
%\setmathfont{Asana-Math.ttf}
% -----------------------------------------------%
% Misc
\author{Κασωτάκη Ε. - Σμαραγδάκης Κ.}
\title{www.math24.gr}
% ---------------- Formation --------------------%


\setlength\topmargin{-0.7cm}
\addtolength\topmargin{-\headheight}
\addtolength\topmargin{-\headsep}
\setlength\textheight{26cm}
\setlength\oddsidemargin{-0.54cm}
\setlength\evensidemargin{-0.54cm}
%\setlength\marginparwidth{1.5in}
\setlength\textwidth{17cm}

\RequirePackage[avantgarde]{quotchap}
\renewcommand\chapterheadstartvskip{\vspace*{0\baselineskip}}
\RequirePackage[calcwidth]{titlesec}
\titleformat{\section}[hang]{\bfseries}
{\Large\thesection}{12pt}{\Large}[{\titlerule[0.9pt]}]
%--------------------------------------------------%
\usepackage{draftwatermark}
\SetWatermarkText{www.math24.gr}
\SetWatermarkLightness{0.9}
\SetWatermarkScale{0.6}
%--------------------------------------------------%


\usepackage{xcolor}
\usepackage{amsthm}
\usepackage{framed}
\usepackage{parskip}

\colorlet{shadecolor}{bluesite!20}

\newtheorem*{orismos}{Ορισμός}

\newenvironment{notation}
  {\begin{shaded}\begin{theorem}}
  {\end{theorem}\end{shaded}}




% Document begins
\begin{document}
\pagestyle{fancy}
\fancyhead{}
\fancyfoot{}
\renewcommand{\headrulewidth}{0pt}
\renewcommand{\footrulewidth}{0pt}


\fancyhead[LO,LE]{
 %\textblockcolor{backsite}
 %\begin{textblock}{5}(-2,-0.55)
  %\rule{0cm}{1cm}
 %\end{textblock}
 \textblockcolor{bluesite}
 \begin{textblock}{5}(-1.5,-0.55)
  \rule{0cm}{1cm}
 \end{textblock}
 %\textblockcolor{bluesite}
 %\begin{textblock}{14}(5.5,-0.55)
 % \rule{0cm}{1cm}
 %\end{textblock}
 \begin{textblock}{0}(-1,-0.25)
 \color{cwhite} \begin{Large}www.math24.gr\end{Large}
 \end{textblock}
\begin{textblock}{0}(16,-0.4)
 \color{cwhite} \includegraphics[height=1.5cm]{math24_logo.png}
 \end{textblock}
\textblockcolor{white}
\begin{textblock}{17}(-1,28)
 \color{backsite} \begin{small}Copyright \textcopyright 2018, Κασωτάκη Ε.(ikasotaki@gmail.com) - Σμαραγδάκης Κ.(kesmarag@gmail.com)\end{small}
 \end{textblock}
}
 
\begin{shaded}
\begin{center}
\huge \textbf{Φυλλάδιο Ασκήσεων}\\
%Πρόσθεση ρητών αριθμών
\end{center} 
\textbf{Μαθηματικά Γ' Γυμνασίου} \hfill \textbf{Ημερομηνία Παράδοσης : \hspace{2em} }
\subsection*{Ονοματεπώνυμο :\hfill  \hspace{5em}}
\end{shaded}
\vspace{2em}
\begin{itemize}
 \item Πρόσθεση, αφαίρεση, πολλαπλασιασμός και διαίρεση πραγματικών αριθμών
 \item Απόλυτη τιμή πραγματικού αριθμού
\end{itemize}
\section*{Θεωρία - Υπενθυμίσεις \hfill \small{}}
\begin{itemize}
 \item \textbf{Φυσικοί αριθμοί} είναι οι αριθμοί $0,1,2,3,4,\cdots ,11,12,\cdots , 997, \cdots$
 \item \textbf{Ακέραιοι αριθμοί} είναι οι φυσικοί και οι αντίθετοί τους, δηλαδή $\cdots ,-3,-2,-1,0,1,2,3,\cdots$ 
 \item \textbf{Ρητοί αριθμοί} είναι οι αριθμοί που μπορούν να πάρουν τη μορφή κλάσματος $\dfrac{μ}{ν}$ 
        όπου $μ,ν$ είναι ακέραιοι αριθμοί με $ν\neq0$\\
        \textbf{π.χ} $\dfrac{5}{3},4$
 \item \textbf{Άρρητοι} είναι οι αριθμοί που δεν είναι ρητοί
 \item \textbf{Η απόλυτη τιμή} ενός αριθμού $α$ εκφράζει την απόσταση του σημείου με τετμημένη $α$ από την 
              την αρχή $Ο$ του άξονα και συμβολίζεται με $|α|$\\
       \textbf{π.χ} $|-7|=7$
\item \textbf{Αντίθετοι αριθμοί} ονομάζονται δύο αριθμοί που είναι ετερόσημοι και έχουν την ίδια απόλυτη τιμή.\\
       \textbf{π.χ} οι αριθμοί $ 3,-3 $ είναι αντίθετοι
\end{itemize}



\section*{Άσκηση 1  \hfill \small{25 μονάδες}}
Να συμπληρώσετε τον παρακάτω πίνακα επιλέγοντας \textbf{"X"} στην κατάλληλη θέση:
\begin{center}
 \begin{tabular}{|c|c|c|c|c|}
\hline                   
       & \textbf{Ακέραιος} &  \textbf{Φυσικός} & \textbf{Ρητός}&  \textbf{Άρρητος}   \\
\hline 
  $8$            &                        &                      &                      &           \\
\hline 
  $-5$            &                        &                       &                     &           \\
\hline 
  $\dfrac{6}{7}$             &                       &                       &                     &           \\
\hline 
  $\sqrt{4}$             &                       &                       &                     &           \\
\hline
  $\sqrt{5}$             &                       &                       &                        &           \\
\hline
  $\sqrt{7}$             &                       &                       &                        &           \\
\hline
  $7,2$             &                       &                       &                        &           \\
\hline
  $\sqrt{100}$             &                       &                      &                       &           \\
\hline
  $6$             &                       &                       &                        &        \\
\hline 
$-5,9$            &                       &                      &                         &   \\
\hline
\end{tabular}
\end{center}







\section*{Άσκηση 2  \hfill \small{25 μονάδες}}
Να συμπληρώσετε τον παρακάτω πίνακα: 
\begin{center}
 \begin{tabular}{|c|c|c|c|c|}
\hline                   
       & \textbf{Αντίθετος} &  \textbf{Αντίστροφος} & \textbf{Απόλυτη τιμή}   \\
\hline 
  $8$            &                        &                      &             \\
\hline 
  $-3$            &                        &                       &            \\
\hline 
  $\dfrac{1}{4}$             &                       &                       &    \\
\hline 
  $-\dfrac{1}{5}$             &                       &                       &         \\
\hline
  $\dfrac{3}{7}$          &                       &                       &          \\
\hline
  $-\dfrac{8}{5}$             &                       &                       &           \\
\hline
  $35$             &                       &                       &                 \\
\hline
  $-71$             &                       &                      &            \\
\hline
  $\dfrac{2}{5}$             &                       &                       &                    \\
\hline 
$-\dfrac{4}{8}$            &                       &                      &                     \\
\hline
\end{tabular}
\end{center}

\section*{Άσκηση 3  \hfill \small{} 25 μονάδες}
Να υπολογίσετε τις παρακάτω παραστάσεις:
\begin{enumerate}[1)]
 \item $3\cdot (\dfrac{1}{3}-\dfrac{2}{3})+4\cdot(\dfrac{7}{2}+\dfrac{6}{12})$
 \item $8\cdot (2-1+4-|-4|)$
 \item $7:\dfrac{8}{3}-\dfrac{1}{2}+\dfrac{3}{4}$
 \item $\dfrac{-\dfrac{1}{2}-\dfrac{2}{3}}{\dfrac{1}{4}+\dfrac{2}{3}}$
 \item $(\dfrac{3}{5}-\dfrac{2}{10}):(\dfrac{7}{3}-\dfrac{10}{6})$
 \item $-3(8-7+\dfrac{1}{2}-\dfrac{1}{3})+|-1|$
 \item $2\cdot |3|+|4|-|-6|-2\cdot \dfrac{1}{4}$
 \item $4(\dfrac{1}{4}+\dfrac{3}{4}-\dfrac{5}{4})+5(\dfrac{6}{5}-\dfrac{18}{10}-\dfrac{16}{20})$
 \item $\dfrac{3+7-2\cdot \dfrac{1}{2}}{5+1+6\cdot \dfrac{1}{3}}$
 \item $5\cdot|-7|-4\cdot|-5|-3\cdot5$
\end{enumerate}



\section*{Άσκηση 4  \hfill \small{} 25 μονάδες}
Αν $α=1$, $β=2$ και $γ+δ=4$, να βρείτε τις τιμές των παρακάτω παραστάσεων:
\begin{enumerate}[i)]
 \item $Α=2α-4β+13$
 \item $Β=6(γ+δ)-|-4|\cdot 5$
 \item $Γ=2δ+2(γ+α)$
 \item $Δ=-3(δ+β)-2(α-δ)-γ$
\end{enumerate}


 









%\let\thefootnote\relax\footnotetext{Επιμέλεια : Κασωτάκη Ειρήνη (ikasotaki@gmail.com),    Σμαραγδάκης Κώστας (kesmarag@gmail.com)} 


\end{document}
